\subsection{CHARACTERIZATION OF AUTOMORPHISMS OF A FINITE-DIMENSIONAL FREE MODULE}

THEOREM 1. Let M be a finite-dimensional free A-module and u an endomorphism of M. The following conditions are equivalent:

\begin{enumerate}
\def\labelenumi{(\alph{enumi})}
\item
  u is bijective;
\item
  \(u\) is right invertible (II, § 1, no. 9, Corollary 1 to Proposition 15);
\item
  u is left invertible (II, § 1, no. 9, Corollary 2 to Proposition 15);
\item
  u is surjective;
\item
  det u is invertible in A.
\end{enumerate}

Let \((e_i)_{1\leqslant i\leqslant n}\) be a basis of M. If \(x_{i} = u(e_{i})\) for \(1\leqslant i\leqslant n\) , then

\[
x _ {1} \wedge x _ {2} \wedge \dots \wedge x _ {n} = (\det u) e _ {1} \wedge e _ {2} \wedge \dots \wedge e _ {n}.
\]

By § 7, no. 9, Theorem 2 a necessary and sufficient condition for the \(x_{i}\) to form a basis of M is that \(\det u\) be an invertible element of A; this proves the equivalence of (a) and (e). Observe that (a) obviously implies each of conditions (b), (c) and (d); it remains to prove that each of conditions (b), (c) and (d) implies (e). Now, if there exists an endomorphism \(v\) of M such that \(v \circ u = 1_{\mathbf{M}}\) or \(u \circ v = 1_{\mathbf{M}}\) , then \((\det v)(\det u) = 1\) and hence \(\det u\) is invertible in A. If \(u\) is surjective, so is \(\bigwedge^{n}(u)\) ( \(\S 7\) , no. 2, Proposition 3), in other words the homothety of ratio det \(u\) in A is surjective, which immediately implies that det \(u\) is invertible.

PROPOSITION 3. Let M be a finite-dimensional free A-module. For every endomorphism u of M, the following conditions are equivalent:

\begin{enumerate}
\def\labelenumi{(\alph{enumi})}
\setcounter{enumi}{5}
\item
  u is injective;
\item
  det u is not a divisor of zero in A.
\end{enumerate}

With the same notation as in the proof of Theorem 1, for u to be injective, it is necessary and sufficient that the \(x_{i}\) be linearly independent. By §7, no. 9, Proposition 12, it is necessary and sufficient for this that the relation \(\lambda x_{1} \wedge x_{2} \wedge \cdots \wedge x_{n} = 0\) (with \(\lambda \in A\) ) imply \(\lambda = 0\) . But this is equivalent to \(\lambda(\det u) = 0\) since \(e_{1} \wedge \cdots \wedge e_{n}\) is a basis of \(\bigwedge^{n}(M)\) ; whence the proposition.

Remark. When A is a field, condition (e) of Theorem 2 is equivalent to condition (g) of Proposition 3, since they both mean that det \(u \neq 0\) . In this case therefore all the conditions of Theorem 1 and Proposition 3 are equivalent (cf.~II, §7, no. 4, Corollary to Proposition 9).

